\chapter{Boltzmann 统计}

\section{Boltzmann 因子}
在第2、3章中，我们处理的都是孤立系统，其中所有允许的态概率都相同，拥有这种简单概率分布的假想系统的集合称为\textbf{微正则系综}。现在我们考虑一个\textit{非孤立系统}，这个系统在某个温度 $T$ 下和热库处于热平衡，我们想求得系统处在某个特定微观态 $s$ 的概率。

假设系统只有一个原子，微观态对应各个能级（含简并态）。如果原子是孤立的，则其能量固定，此时具有该能量的所有微观态等可能；但现在其与热库自由交换能量，故态的概率取决于能量。

考虑两个态 $s_1,s_2$，能量分别为 $E(s_1),E(s_2)$，概率分别为 $\mathcal{P}(s_1),\mathcal{P}(s_2)$，对应热库的微观态重数（或熵）分别为 $\Omega_R(s_1),\Omega_R(s_2)$（或 $S_R(s_1),S_R(s_2)$），则:
\begin{equation}
    \frac{\mathcal{P}(s_1)}{\mathcal{P}(s_2)}
    =\frac{\Omega_R(s_1)}{\Omega_R(s_2)}
    =\frac{e^{S_R(s_1)/k}}{e^{S_R(s_2)/k}}
    =\exp\!\left(\frac{S_R(s_1)-S_R(s_2)}{k}\right)
\end{equation}

原子相较于热库很小，故：
\begin{equation}\label{eq:deltaS}
    S_R(s_2)-S_R(s_1)=\mathrm{d}S_R
    =\frac{1}{T}\left(\mathrm{d}U_R+P\,\mathrm{d}V_R-\mu\,\mathrm{d}N_R\right)
\end{equation}
取 $\mathrm{d}N_R=0$ 且 $P\,\mathrm{d}V_R\ll \mathrm{d}U_R$，得：
\begin{equation}
    S_R(s_2)-S_R(s_1)
    \approx \frac{1}{T}\,\mathrm{d}U_R
    =\frac{1}{T}\bigl[U_R(s_2)-U_R(s_1)\bigr]
    =\frac{-\bigl[E(s_2)-E(s_1)\bigr]}{T}
\end{equation}
因此：
\begin{equation}
    \frac{\mathcal{P}(s_2)}{\mathcal{P}(s_1)}
    =\exp\!\left(-\frac{E(s_2)-E(s_1)}{kT}\right)
    =\frac{e^{-E(s_2)/kT}}{e^{-E(s_1)/kT}}
    \implies \mathcal{P}(s)\propto e^{-E(s)/kT}
\end{equation}

\begin{definition}[Boltzmann 因子]
    \begin{equation}
        \text{Boltzmann 因子} \equiv e^{-E(s)/kT} = e^{-E(s)\beta},\qquad \text{其中} \beta = (kT)^{-1}
    \end{equation}
\end{definition}

进一步作归一化，得到配分函数：
\begin{equation}
    \mathcal{P}(s)=\frac{1}{Z}e^{-\beta E(s)} \implies Z=\sum_s e^{-\beta E(s)}
\end{equation}

\begin{definition}[配分函数]
    \begin{equation}
        Z \equiv \sum_s e^{-E(s)/kT}
    \end{equation}
\end{definition}


\section{平均值}
我们现在知道了当系统与一个温度为 $T$ 的热库平衡时，对于任意的微观态 $s$，其概率由 \textbf{Boltzmann 分布} 给出：
\begin{definition}[Boltzmann 分布]
    \begin{equation}
        \mathcal{P}(s) = \frac{1}{Z} e^{\beta E(s)}
    \end{equation}
\end{definition}
所有这种概率服从 Boltzmann 分布的假想系统的集合称为\textbf{正则系综}。

如果我们不关心系统在每个微观态上的概率的取值，而是对系统某个性质的\textit{平均值}感兴趣，该如何计算？ 为此，设系统有 $N$ 个粒子，位于 $s$ 态具有 $N(s)$ 个，则：
\begin{equation}
    \overline E
    =\frac{1}{N}\sum_s E(s)N(s)
    =\sum_s E(s)\mathcal{P}(s)
    =\frac{1}{Z}\sum_ E(s)e^{-\beta E(s)}
\end{equation}
注意到：
\begin{equation}
    \frac{\partial Z}{\partial \beta}=-\sum_s E(s)e^{-\beta E(s)}
    \implies \overline E=-\frac{1}{Z}\frac{\partial Z}{\partial \beta}
    =-\frac{\partial}{\partial \beta}\ln Z
\end{equation}

若系统有一系列相同粒子，则可以用 $U=N\overline E$ 代表系统平均能量。但系统的能量存在涨落，考虑标准差 $\sigma_E $：
\begin{equation}
    \overline{E^2}=\frac{1}{Z}\sum_s E(s)^2e^{-\beta E(s)}=\frac{1}{Z}\frac{\partial^2 Z}{\partial \beta^2}
\end{equation}
故：
\begin{equation}  
    \begin{split}
        \sigma_E^2 &=\overline{E^2}-\overline E^2\\
        &=\frac{1}{Z}\frac{\partial^2 Z}{\partial \beta^2}-\frac{1}{Z^2}\left(\frac{\partial Z}{\partial \beta}\right)^2\\
        &=\frac{\partial}{\partial \beta}\left[\frac{1}{Z}\frac{\partial Z}{\partial \beta}\right]
        =-\frac{\partial \overline E}{\partial \beta}
        = -\frac{\partial \overline E}{\partial T}\frac{\partial T}{\partial \beta}\\
        &= kT^2 C        
    \end{split}
\end{equation}
因此：
\begin{corollary}[正则系综的能量涨落]
    \begin{equation}
        \sigma_E=kT\sqrt{C}
    \end{equation}
\end{corollary}

容易想象 $E\propto N,\ C\propto N$，于是
\begin{equation}
    \frac{\sigma_E}{\overline E}\propto \frac{1}{\sqrt{N}}
    \implies
    N\text{ 足够大时相对涨落}\to 0
\end{equation}
所以对于宏观的热力学系统，能量涨落并不显著，平均能量可以很好地反映系统的总能量。

\subsection{双状态顺磁体}
回忆理想双状态顺磁体中每一个偶极子只有两个可能的态：$-\mu B$ 和 $\mu B$，所以每个偶极子的配分函数是：
\begin{equation}
    Z=e^{\beta\mu B}+e^{-\beta\mu B}=2\cosh(\beta\mu B)
\end{equation}
所以：
\begin{align}
    &\overline E=-\frac{1}{Z}\frac{\partial Z}{\partial \beta}
    =-\mu B\tanh(\beta\mu B) \implies\\
    &U= N \overline E = -N\mu B\tanh(\beta\mu B)\\
    &\overline\mu_z=\sum_s \mu(s)\mathcal{P}(s)=\mu\tanh(\beta\mu B),\\
    &M=N\overline\mu_z=N\mu\tanh(\beta\mu B)
\end{align}
与 3.3 节中的推导结果一致。

\subsection{Einstein 固体}
考虑在温度 $T$ 下 $N$ 个相同谐振子，每个谐振子允许的能量为 $0,hf,2hf,\dots$，则：
\begin{align}
&Z=\sum_{n=0}^{\infty}e^{-\beta nhf}=\frac{1}{1-e^{-\beta hf}}\\
&\overline E=-\frac{1}{Z}\frac{\partial Z}{\partial \beta} =\frac{hf\,e^{-\beta hf}}{1-e^{-\beta hf}} =\frac{hf}{e^{\beta hf}-1}\\
&U=\frac{Nhf}{e^{\beta hf}-1}
\end{align}
与习题 3.25 的结果一致。

\subsection{双原子分子转动}
假设分子是孤立的，其转动能量为
\begin{equation}
    E(j)=j(j+1)\,\varepsilon
\end{equation}
其中能级 $j$ 对应 $2j+1$ 个简并态。对 $j$ 求和写出配分函数：
\begin{equation}
Z_{\mathrm{rot}}=\sum_{j=0}^{\infty}(2j+1)e^{-\beta j(j+1)\varepsilon}
\end{equation}
在高温极限 $kT\gg\varepsilon$ 下，可以用积分近似求和：
\begin{align}
&Z_{\mathrm{rot}}
\approx \int_0^{\infty}(2j+1)e^{-\beta j(j+1)\varepsilon}\,\mathrm{d}j
=\frac{1}{\beta\varepsilon}=\frac{kT}{\varepsilon}\\
&\overline E_{\mathrm{rot}}
=-\frac{1}{Z_{\mathrm{rot}}}\frac{\partial Z_{\mathrm{rot}}}{\partial \beta}
=\frac{1}{\beta}=kT
\end{align}
这与能量均分定理给出的结果一致（2 个自由度）。若分子的两个原子相同，相当于态数目减半，则
\begin{equation}
Z_{\mathrm{rot}}\approx \frac{kT}{2\varepsilon}
\end{equation}
但 $\overline E_{\mathrm{rot}}$ 不变。

\section{能量均分定理}
在 1.3 节中我们讨论过能量均分定理，现在我们可以给出证明。
\begin{theorem}[能量均分定理]
    对于能量与``自由度''呈二次方关系（即有 $E(q) = cq^2$ 形式）的系统，其能量满足：
    \begin{equation}
        \overline{E} = \frac{1}{2}kT
    \end{equation}
\end{theorem}

我们先假设这个系统离散，每个 $q$ 对应一个独立状态，状态之间间隔为 $\Delta q$，再令 $\Delta q \to 0$，则：
\begin{equation}
    Z
    =\sum_q e^{-\beta E(q)}
    =\sum_q e^{-\beta cq^2}
    \xlongequal{\Delta q\to 0} \frac{1}{\Delta q}\int_{-\infty}^{\infty}e^{-\beta cq^2}\,\mathrm{d}q
\end{equation}
令 $x=\sqrt{\beta c}\,q$，则：
\begin{equation}
Z
=\frac{1}{\Delta q}\frac{1}{\sqrt{\beta c}}\int_{-\infty}^{\infty}e^{-x^2}\,\mathrm{d}x
=\frac{1}{\Delta q}\sqrt{\frac{\pi}{\beta c}}
\end{equation}
于是
\begin{equation}
\overline E=-\frac{1}{Z}\frac{\partial Z}{\partial \beta}
=\frac{1}{2}\frac{1}{\beta}=\frac{1}{2}kT
\end{equation}

\begin{remark}
    能量均分定理只适用于经典系统且能级间距 $\ll kT$，此时 Gauss 曲线才比较好地近似求和；对于量子力学系统，能量均分定理是不成立的。
\end{remark}

\section{Maxwell 速率分布}
Boltzmann 因子的另一个应用是计算理想气体分子运动速率的概率。严格来说，分子速率\textit{是}某个定值 $v$ 的概率是\textit{零}，因为速率可以连续变化。因此，这里需要用分布函数重新定义概率：
\begin{definition}[分布函数与概率]
    定义速率分布函数 $\mathcal{D}(v)$，它在任意一个区间上的积分表示分子速率落在这个区间内的概率：
    \begin{align}
        & \mathcal{P}(v_1<v<v_2) = \int_{v_1}^{v_2}\mathcal{D}(v)\,\mathrm{d}v\\
        & \mathcal{P}(v_1<v<v_1+\,\mathrm{d}v) = \mathcal{D}(v)\,\mathrm{d}v
    \end{align}
\end{definition}

可以想象，$D(v)$ 正比于（一个分子拥有 $v$ 的概率）$\times$（对应 $v$ 的态数目）。前者正比于 Boltzmann 因子，后者则正比于速度矢量空间中一个半径为 $v$ 的球的表面积：
\begin{equation}
    \mathcal{D}(v)\propto e^{-\frac{1}{2}\beta mv^2}\cdot4\pi v^2
\end{equation}
设比例系数为 $C$，将概率归一化：
\begin{equation}
    1=\int_0^{\infty}\mathcal{D}(v)\,\mathrm{d}v
    =4\pi C\int_0^{\infty}v^2e^{-\frac{1}{2}\beta mv^2}\,\mathrm{d}v
    = C\left(\frac{2\pi}{\beta m}\right)^{3/2}
\end{equation}
故：
\begin{equation}
    C=\left(\frac{\beta m}{2\pi}\right)^{3/2}=\left(\frac{m}{2\pi kT}\right)^{3/2}
\end{equation}
最终我们得到理想气体分子速率的 \textbf{Maxwell 分布}：
\begin{corollary}[Maxwell 分布]
    理想气体分子的速率分布近似由分布函数 $\mathcal{D}(v)$ 确定：
    \begin{equation}
        \mathcal{D}(v) = \left(\frac{m}{2\pi kT}\right)^{3/2}4\pi v^2e^{-\frac{mv^2}{2kT}}
    \end{equation}
\end{corollary}

进一步我们可以比较三种描述理想气体分子速率的量：
\begin{equation}
v_{\max}=\sqrt{\frac{2kT}{m}},\qquad
\overline v=\int_0^{\infty}v\mathcal{D}(v)\,\mathrm{d}v=\sqrt{\frac{8kT}{\pi m}},\qquad
v_{\mathrm{rms}}=\sqrt{\frac{3kT}{m}}
\end{equation}

\begin{figure}[htbp]
  \centering
  \resizebox{0.9\linewidth}{!}{%
    \begin{tikzpicture}[
        x=1cm,y=1cm,
        >=Stealth,
        every node/.style={font=\small}
        ]
        % ---------- parameters (tune if you want) ----------
        \def\xmax{10.2}      % x-axis length
        \def\ymax{3.6}       % y-axis length
        \def\vmaxpos{3.5}    % v_max tick position
        \def\vbarpos{4.0}    % v_bar tick position
        \def\vrmspos{4.7}    % v_rms tick position

        % Maxwell-like curve: y = A * x^2 * exp(-x^2/a)
        % Choose a = vmax^2 so peak is near v_max tick.
        \def\a{12.25}        % = (3.5)^2
        \def\A{0.7}         % vertical scale

        % ---------- axes ----------
        \draw[-{Stealth[length=4mm,width=3.2mm]}, line width=0.9pt] (0,0) -- (\xmax,0);
        \draw[-{Stealth[length=4mm,width=3.2mm]}, line width=0.9pt] (0,0) -- (0,\ymax);

        \node[anchor=east] at (-0.25,\ymax-0.1) {$\mathcal{D}(v)$};
        \node[anchor=west] at (\xmax+0.1,0.02) {$v$};

        % ---------- curve ----------
        \draw[line width=2.2pt]
        plot[smooth,domain=0:\xmax-0.2,samples=250]
        (\x,{ \A*(\x*\x)*exp(-(\x*\x)/\a) });

        % ---------- annotations ----------
        \node[anchor=east] (par) at (-0.1,1.25) {抛物线};
        \draw[line width=0.9pt] (par.south) -- (0.65,0.40);

        \node[anchor=west] (die) at (7.2,1.35) {指数衰减};
        \draw[line width=0.9pt] (die.south) -- (7.65,0.55);

        % ---------- ticks + labels (with the little “brackets”) ----------

        \draw[line width=0.8pt] (\vmaxpos,0.22) -- (\vmaxpos,-0.14);
        \draw[line width=0.8pt] (\vbarpos,0.22) -- (\vbarpos,-0.38);
        \draw[line width=0.8pt] (\vrmspos,0.22) -- (\vrmspos,-0.14);

        % small “bracket” connectors like the original (left and right)
        \draw[line width=0.8pt] (\vmaxpos,-0.14) -- (\vmaxpos-0.35,-0.30);
        \draw[line width=0.8pt] (\vrmspos,-0.14) -- (\vrmspos+0.35,-0.30);

        % labels
        \node[below] at (\vmaxpos-0.55,-0.38) {$v_{\max}$};
        \node[below] at (\vbarpos,-0.38) {$\overline v$};
        \node[below] at (\vrmspos+0.55,-0.38) {$v_{\mathrm{rms}}$};

    \end{tikzpicture}%
  }
  \caption{Maxwell 速率分布示意图}
\end{figure}

\section{配分函数与自由能}
回忆 $F\equiv U-Ts$ 与偏导关系，有：
\begin{equation}
    \left(\frac{\partial F}{\partial T}\right)_{V,N}=-S
    \implies
    \left(\frac{\partial F}{\partial T}\right)_{V,N}=\frac{F-U}{T}
\end{equation}
另一方面，设 $\widetilde{F}=-kT\ln Z$，则：
\begin{equation}
    \left(\frac{\partial \widetilde{F}}{\partial T}\right)_{V,N}
    =\frac{\partial}{\partial T}\bigl(-kT\ln Z\bigr)
    =-k\ln Z-kT\frac{\partial}{\partial T}\ln Z
\end{equation}
又：
\begin{equation}
    \frac{\partial}{\partial T}\ln Z
    =\frac{\partial \beta}{\partial T}\frac{\partial}{\partial \beta}\ln Z
    =\left(-\frac{1}{kT^2}\right)(-U)
    =\frac{U}{kT^2}
\end{equation}
代回得：
\begin{equation}
    \left(\frac{\partial \widetilde{F}}{\partial T}\right)_{V,N}
    =-k\ln Z-\frac{U}{T}=\frac{\widetilde{F}-U}{T}
\end{equation}
我们发现 $\widetilde{F}$ 和 $F$ 满足同一个微分方程。同时在 $T=0$ 处，$F = U = U_0$，而配分函数为 $Z = e^{U_0/kT}$，因此 $F(0) = \widetilde{F}(0)$。所以对所有的 $T$，$F=\widetilde{F}$。我们得到了配分函数和自由能的重要联系：
\begin{corollary}[配分函数和自由能]
    \begin{equation}
        F = -kT \ln Z
    \end{equation}
\end{corollary}
从 $F$ 出发，我们可以计算系统的其他热力学性质：
\begin{equation}
    S=-\left(\frac{\partial F}{\partial T}\right)_{V,N},\qquad
    P=-\left(\frac{\partial F}{\partial V}\right)_{T,N},\qquad
    \mu=\left(\frac{\partial F}{\partial N}\right)_{T,V}    
\end{equation}

\section{组合系统的配分函数}
考虑 $N$ 个粒子的系统，总能量为 $\sum_i E_i$，且粒子间无作用、可区分，则：
\begin{equation}
    \begin{split}
        Z_{\mathrm{tot}}
        &=\sum_s e^{-\beta\sum_i E_i(s)}
        =\sum_{s_1}\sum_{s_2}\cdots\sum_{s_N}\prod_i e^{-\beta E_i(s_i)}\\
        &=\prod_i\left[\sum_{s_i}e^{-\beta E_i(s_i)}\right]
        =\prod_i Z_i
        =Z_1Z_2\cdots Z_N
    \end{split}
\end{equation}
其中 $s_i$ 为第 $i$ 个粒子的微观态。若粒子间不可区分，则：
\begin{equation}
    Z_{\mathrm{tot}}=\frac{1}{N!}Z_1Z_2\cdots Z_N
\end{equation}

\section{再看理想气体}
考虑包含 $N$ 个相同分子的理想气体：
\begin{equation}
    Z=\frac{1}{N!}Z_1^N
\end{equation}
单个分子的能量由平动能 $E_{\mathrm{tr}}$ 和内部能 $E_{\mathrm{int}}$ 组成：
\begin{equation}
    e^{-\beta E(s)}=e^{-\beta E_{\mathrm{tr}}(s)}\,e^{-\beta E_{\mathrm{int}}(s)}
    \implies
    Z_1=Z_{\mathrm{tr}}Z_{\mathrm{int}}
\end{equation}

为计算 $Z_{\mathrm{tr}}$，想象分子是被限制在盒子内的波函数：先考虑 1D 情况，由于波函数在两端为 0，允许的驻波波长为：
\begin{equation}
    \lambda_n=\frac{2L}{n},\qquad n=1,2,\dots
\end{equation}
于是：
\begin{equation}
    p_n=\frac{h}{\lambda_n}=\frac{hn}{2L} \implies E_n=\frac{p_n^2}{2m}=\frac{h^2n^2}{8mL^2}
\end{equation}
写出配分函数，对于充分大的盒子，将求和近似为积分：
\begin{equation}
    Z_{1\mathrm{D}}=\sum_n e^{-\beta E_n}
    \approx\int_0^{\infty}\exp\!\left(-\frac{h^2n^2}{8mL^2kT}\right)\,\mathrm{d}n
    =\frac{\sqrt{2\pi mkT}}{h}L
    \equiv\frac{L}{\ell_Q}
\end{equation}
其中 $\ell_Q=\frac{h}{\sqrt{2\pi mkT}}$ 是量子长度。推广到 3D：
\begin{equation}
    E_{\mathrm{tr}}=\frac{p_x^2}{2m}+\frac{p_y^2}{2m}+\frac{p_z^2}{2m} \implies Z_{\mathrm{tr}}=\sum_s e^{E_{\mathrm{tr}}/kT} = \prod_\mathrm{cyc} \sum_{n_x} e^{E_x/kT} =\left(\frac{L}{\ell_Q}\right)^3 \equiv \frac{V}{v_Q}
\end{equation}
其中 $V$ 是盒子体积，$v_Q=\ell_Q^3=\left(\frac{h}{\sqrt{2\pi mkT}}\right)^3$ 是量子体积。

另一种方式是从经典的假设出发，对每个有可能分配位置和动量，并引入 $1/h^3$ 消除量纲：
\begin{equation}
    \begin{split}
        Z_{\mathrm{tr}}
        &=\frac{1}{h^3}\int \mathrm{d}^3r\,\mathrm{d}^3p\,e^{-E_{\mathrm{tr}}/kT}\\
        &=\frac{1}{h^3}\int \mathrm{d}^3r\int \mathrm{d}^3p\exp\!\left(-\frac{p^2}{2mkT}\right)\\
        &=\frac{V}{h^3}\left(\sqrt{2\pi mkT}\right)^3
        \equiv \frac{V}{v_Q}  
    \end{split}
\end{equation}

总之：
\begin{equation}
    Z_1=Z_{\mathrm{tr}}Z_{\mathrm{int}}=\frac{V}{v_Q}Z_{\mathrm{int}}
    \implies
    Z=\frac{1}{N!}\left(\frac{VZ_{\mathrm{int}}}{v_Q}\right)^N
\end{equation}
利用 Stirling 近似展开：
\begin{equation}
    \ln Z
    =N\ln\left(\frac{VZ_{\mathrm{int}}}{v_Q}\right)-\ln N!
    = N\left(\ln V+\ln Z_{\mathrm{int}}-\ln N-\ln v_Q+1\right)
\end{equation}
我们可以计算出理想气体的一系列热力学性质：
\begin{corollary}[理想气体的热力学性质]
    \begin{align}
        &U=-\frac{\partial}{\partial \beta}\ln Z =-N\frac{\partial}{\partial \beta}\ln Z_{\mathrm{int}}+N\frac{1}{v_Q}\frac{\partial v_Q}{\partial \beta}= N\overline E_{\mathrm{int}}+\frac{3}{2}NkT=U_{\mathrm{int}}+\frac{3}{2}NkT\\
        &C_V=\frac{\partial U}{\partial T}=\frac{\partial U_{\mathrm{int}}}{\partial T}+\frac{3}{2}Nk\\
        &F=-kT\ln Z=-NkT\left(\ln V -\ln N-\ln v_Q+1\right)+  F_{\mathrm{int}}\\
        &P=-\left(\frac{\partial F}{\partial V}\right)_{T,N}=\frac{NkT}{V}\\
        &S=-\left(\frac{\partial F}{\partial T}\right)_{V,N}=Nk\left(\ln\frac{V}{Nv_Q}+\frac{5}{2}\right)-\frac{\partial F_{\mathrm{int}}}{\partial T}\\
        &\mu=\left(\frac{\partial F}{\partial N}\right)_{T,V}= -kT\ln\frac{V}{Nv_Q}-kT\ln Z_{\mathrm{int}} =-kT\ln\left(\frac{VZ_{\mathrm{int}}}{Nv_Q}\right)\\
        &G=F+PV=-NkT\left[\ln\left(\frac{VZ_{\mathrm{int}}}{Nv_Q}\right)+1\right]+NkT=-NkT\ln\left(\frac{VZ_{\mathrm{int}}}{Nv_Q}\right)=N\mu
    \end{align}
\end{corollary}